%% abstract
\begin{abstract}
I introduce \emph{antiTranscription}, an end-to-end computer-vision pipeline for
converting a casual smartphone photo of printed sheet music into a playable MIDI
file. My system consists of five learned and traditional steps: (i) a
homography-based perspective rectification using a reliable page-corner cascade;
(ii) Hough-transform staff-line detection coupled with a five-line clustering
algorithm; (iii) morphological staff removal followed by connected-component
symbol segmentation; (iv) a trained convolutional neural network for glyph
recognition, on a $35{,}933$-example crop set derived from the PrIMuS
corpus~\cite{calvo2018primus}; and (v) geometric pitch reading and MIDI
synthesis. I find that the geometry steps are extremely reliable on five real
phone photos (rectification passes on all five with slant below $0.22^{\circ}$;
staff detection on all pages), while my main contribution shows a large
\emph{clean-to-real} domain gap: the classifier that achieves $58.9\%$ on clean
glyphs falls to $22.7\%$ on real-cropped glyphs, a drop of $36.2$ points,
primarily concentrated on delicate symbols like hollow half notes (recognized
with $0\%$ accuracy). This motivates my approach of using a geometry-grounded,
segmentation-free notehead-localization and pitch reader that I show improves
exact pitch accuracy from $14\%$ to $\mathbf{86\%}$ on a test song ($89\%$ within
one step). OMR using geometric reasoning is much more promising for phone images
than an appearance-based classifier trained on clean data.
\end{abstract}
